\chapter{Moments of \texorpdfstring{$L$}{L}-functions}\label{ch:oc-moments}\label{ch:moments}
The Keating--Snaith conjecture and the recipe that turns a random-matrix moment into a prediction for $\zeta$. What is proved here is the mean square, with Ingham's secondary term, the lower bounds of the conjectured order, and the arithmetic factor $a_k$ read from the random Euler product of Section~\ref{ch:bohrjessen}. The matrix moments are Volume~8. Moments of $\zeta$ on the critical line past the fourth stay conjectural.

The approximate functional equation is not proved in this chapter. Every use of it is marked. The proved mean square does not use it: its secondary term comes from two poles of a zeta factor, as the text says. The lower bounds of the conjectured order are proved from short Dirichlet polynomials and Hölder, without an approximate functional equation. The recipe in the conjecture section writes $\zeta(\frac12+it)^2$ as a sum of length $t/2\pi$ plus a dual sum of the same length; that sentence is a heuristic, and the exercise that uses it says so. The Gonek--Hughes--Keating factorisation of $\zeta$ into an Euler product and a product over zeros is quoted as their theorem about an approximation, and the independence of the two factors is marked as a conjecture. No unmarked line of the recipe treats the approximate functional equation as proved.

\begin{center}\itshape
The mean of $|\zeta|^2$ was known in 1918, the mean of $|\zeta|^4$ in 1926.\\
For the mean of $|\zeta|^6$ the world waited seventy-two years for a conjecture,\\
and the number it predicted was $42$.
\end{center}

Selberg's theorem describes the typical size of $|\zeta(\frac12+it)|$: it is $e^{\pm\sqrt{\frac12\log\log T}\,Z}$ with $Z$ Gaussian. The \emph{moments}
\[
\mathcal M_k(T)=\frac1T\int_T^{2T}|\zeta(\tfrac12+it)|^{2k}\dd t
\]
ask a different question, about the atypical: how large is $|\zeta|$ on the rare occasions when it is large? These averages have been the central problem of the theory of $\zeta$ for a century, partly because they control everything else. A good bound for $\mathcal M_k$ with $k$ large gives a bound for the maximum of $|\zeta|$ (the Lindelöf hypothesis is equivalent to $\mathcal M_k(T)=T^{o(1)}$ for every $k$), and bounds for $\mathcal M_2$ are behind many results on primes in short intervals and zero density.

The first two moments were found early. In 1918 G.~H.~Hardy and J.~E.~Littlewood proved that $\mathcal M_1(T)\sim\log T$. In 1926 Albert Ingham proved that $\mathcal M_2(T)\sim\frac1{2\pi^2}(\log T)^4$. Then nothing, for a very long time: no proof for $k=3$, and not even a convincing guess for the constant. A natural conjecture had the shape $\mathcal M_k(T)\sim c_k(\log T)^{k^2}$, and the constant was written as $c_k=a_k\frac{g_k}{(k^2)!}$ with an explicit arithmetic factor $a_k$ (a product over primes) and an unknown number $g_k$, with $g_1=1$ and $g_2=2$. In 1998 Brian Conrey and Amit Ghosh, after a long number-theoretic computation, conjectured $g_3=42$. Readers of Douglas Adams will recognise the number. Then came the surprise. Jon Keating and Nina Snaith, working with random unitary matrices rather than with $\zeta$, found a formula that gives $g_k$ for every $k$: $1,2,42,24024,\dots$ (Volume~11, Theorem~11.10). It agreed with Conrey and Ghosh, and with the value $g_4=24024$ that Conrey and Steve Gonek had just conjectured by yet another method.

This chapter proves what can be proved with the tools at hand: the mean square, with Ingham's lower-order term; lower bounds $\mathcal M_k(T)\gg(\log T)^{k^2}$ of the conjectured order, for every integer $k$; and the computation, from the random Euler product of Section~\ref{ch:bohrjessen}, of the arithmetic factor $a_k$. The rest is the conjecture, and the numerical evidence for it, which is subtle: at every height a computer can reach, the moments are still far from their conjectured leading term.

\begin{reminder}[from Volumes~4, 5, 11 and~13]
\emph{Near $s=1$} (Volume~4, Proposition~10.6). $\zeta(s)=\frac1{s-1}+\gamma+O(s-1)$, and $\frac1{\sigma-1}\le\zeta(\sigma)\le\frac1{\sigma-1}+1$ for $\sigma>1$.
\emph{Mertens' product} (Volume~5, Theorem~1.6). $\prod_{p\le y}(1-\frac1p)=\frac{e^{-\gamma}}{\log y}(1+O(\frac1{\log y}))$.
\emph{$\zeta$ as a short sum} (Volume~13, Lemma~1.5), as recalled in Section~\ref{ch:bohrjessen}.
\emph{Random unitary matrices} (Volume~11, Theorems~11.10 and~11.11). If $U$ is a Haar unitary matrix of size $N$ and $m\ge0$ is an integer, $\mathbb E|\det(I-U)|^{2m}\sim\frac{G(m+1)^2}{G(2m+1)}N^{m^2}$, where $G$ is Barnes' function, $G(m+1)=0!\,1!\cdots(m-1)!$; and $\log|\det(I-U)|/\sqrt{\frac12\log N}\to\mathcal N(0,1)$.
From Section~\ref{ch:selbergclt}: the mean value Lemma~\ref{lem:mvt2} and the twisted second moment, Theorem~\ref{thm:twisted}.
\end{reminder}

\section{The mean square}

\begin{theorem}[Hardy and Littlewood, 1918; Ingham, 1926]\label{thm:ingham}
As $T\to\infty$,
\[
\int_T^{2T}|\zeta(\tfrac12+it)|^2\dd t=\int_T^{2T}\Bigl(\log\frac t{2\pi}+2\gamma\Bigr)\dd t+O(T^{0.6}).
\]
In particular $\mathcal M_1(T)=\log T+O(1)$.
\end{theorem}
\begin{proof}
Theorem~\ref{thm:twisted} with $h=k=1$, whose constants do not depend on $\sigma\in\left]\frac12,0.6\right]$, gives for those $\sigma$
\[
\int_T^{2T}|\zeta(\sigma+it)|^2\dd t=\zeta(2\sigma)T+\zeta(2-2\sigma)\int_T^{2T}\Bigl(\frac t{2\pi}\Bigr)^{1-2\sigma}\dd t+O\bigl(T^{1/2}(\log T)^3+T^{0.6}\bigr).
\]
Let $\varepsilon=2\sigma-1\to0^+$, with $T$ fixed. The left side tends to $\int_T^{2T}|\zeta(\frac12+it)|^2\dd t$, because $\zeta$ is continuous on the compact set $[\frac12,0.6]\times[T,2T]$. On the right, $\zeta(1+\varepsilon)T=\frac T\varepsilon+\gamma T+O(\varepsilon T)$, and with $\ell(t)=\log\frac t{2\pi}$, $(\frac t{2\pi})^{-\varepsilon}=1-\varepsilon\ell(t)+O(\varepsilon^2\log^2T)$, so
\[
\begin{aligned}
\zeta(1-\varepsilon)\int_T^{2T}\Bigl(\frac t{2\pi}\Bigr)^{-\varepsilon}\dd t
&=\Bigl(-\frac1\varepsilon+\gamma+O(\varepsilon)\Bigr)\Bigl(T-\varepsilon\int_T^{2T}\ell+O(\varepsilon^2T\log^2T)\Bigr)\\
&=-\frac T\varepsilon+\int_T^{2T}\ell+\gamma T+O(\varepsilon T\log^2T).
\end{aligned}
\]
The poles in $\varepsilon$ cancel, and letting $\varepsilon\to0$ gives the formula. Since $\int_T^{2T}\ell=T\log T+O(T)$, $\mathcal M_1(T)=\log T+O(1)$.
\end{proof}

Two features of the proof deserve a comment. The main term comes from two poles that collide at $\sigma=\frac12$: one from the diagonal $\sum_nn^{-2\sigma}=\zeta(2\sigma)$, the other from the functional equation. Off the line they are separate, and the second is negligible (it was $o(T)$ in Proposition~\ref{prop:mollmeansq}); on the line they merge into the logarithm. And the size $\log T$ has a plain meaning: it is the number $\sum_{n\le T}\frac1n$ of ``independent'' terms $n^{-1/2-it}$ that a sum of length $T$ can hold.

\section{Why the exponent is \texorpdfstring{$k^2$}{k squared}}

Here is a heuristic, which is also the first hint that random matrices know something. Suppose that Selberg's theorem were exact: $\log|\zeta(\frac12+it)|$ Gaussian with mean $0$ and variance $\frac12\log\log T$. Then, since $\mathbb Ee^{\lambda Z}=e^{\lambda^2v/2}$ for $Z\sim\mathcal N(0,v)$,
\[
\mathbb E_T|\zeta(\tfrac12+it)|^{2k}=\mathbb E_Te^{2k\log|\zeta|}=e^{\frac12(2k)^2\cdot\frac12\log\log T}=(\log T)^{k^2}.
\]
The exponent $k^2$ is exactly the conjectured one. The same computation for a Haar unitary matrix, with the variance $\frac12\log N$ of Volume~11, Theorem~11.11, gives $N^{k^2}$, the exponent of Volume~11, Theorem~11.10. In both cases the constant is wrong, and must be: the moment is carried by the values where $\log|\zeta|\approx k\log\log T$, which have probability about $(\log T)^{-k^2}$ (Exercise~\ref{ex:lognormal}), deep in the tail, where a central limit theorem says nothing. Figure~\ref{fig:momentshare} shows how extreme this is already for small $k$.

\begin{figure}[ht]
\centering\fig{ch08-share}
\caption{Where the moments come from, at $T=10^6$ ($200\,000$ random points $t\in[T,2T]$, Riemann--Siegel formula). The $1\%$ of points where $|\zeta(\frac12+it)|$ is largest carry $31\%$ of $\int|\zeta|^2$, $78\%$ of $\int|\zeta|^4$ and $95\%$ of $\int|\zeta|^6$; the top $0.1\%$ carry $7\%$, $33\%$ and $62\%$. The dashed diagonal is what equal contributions would give.}
\label{fig:momentshare}
\end{figure}

\begin{think}
If the top $0.1\%$ of points carry most of $\int|\zeta|^6$, how many points must a Monte Carlo estimate of $\mathcal M_3(T)$ use to reach $10\%$ accuracy? What does this say about the numerical evidence for $g_3=42$?
\end{think}

\section{Lower bounds of the right order}

An upper bound $\mathcal M_k(T)\ll(\log T)^{k^2}$ is out of reach for $k\ge3$, even on the Riemann hypothesis it was proved only in 2013. Lower bounds are easier, and a short argument of Zeév Rudnick and Kannan Soundararajan gives them with the right exponent. The idea is to compare $\zeta$ with a short Dirichlet polynomial $A$ that imitates it: if $\zeta$ correlates with $A$, and $A$ is not too large, then $\zeta$ must be large. Hölder's inequality does the rest.

We need the size of a divisor sum. For an integer $k\ge1$, $d_k(n)$ is the number of ordered $k$-tuples of positive integers with product $n$, so that $\zeta(s)^k=\sum_nd_k(n)n^{-s}$ for $\sigma>1$, and $d_k(p^j)=\binom{j+k-1}{k-1}$.

\begin{lemma}\label{lem:dk2}
For each integer $k\ge1$ there are $0<c_k\le C_k$ such that for $x\ge3$,
\[
c_k(\log x)^{k^2}\le\sum_{n\le x}\frac{d_k(n)^2}n\le C_k(\log x)^{k^2}.
\]
\end{lemma}
\begin{proof}
For $\sigma>1$ let $F(\sigma)=\sum_nd_k(n)^2n^{-\sigma}=\prod_ph(p^{-\sigma})$, with $h(u)=(1-u)^{-k^2}\cdot h_0(u)$ and $h_0(u)=(1-u)^{k^2}\sum_{j\ge0}\binom{j+k-1}{k-1}^2u^j$. The power series $h_0$ has radius $1$, $h_0(0)=1$, and its coefficient of $u$ is $k^2-k^2=0$; so $|h_0(u)-1|\le M_ku^2$ for $0\le u\le\frac12$, and $h_0(u)>0$ there. Hence $g(\sigma)=\prod_ph_0(p^{-\sigma})$ converges and lies between two positive constants for $1\le\sigma\le2$: the factors with $M_kp^{-2}\le\frac12$ lie in $[e^{-2M_kp^{-2}},e^{M_kp^{-2}}]$, and the finitely many others are continuous and positive in $\sigma$. Since $F(\sigma)=\zeta(\sigma)^{k^2}g(\sigma)$ and $\zeta(\sigma)\asymp\frac1{\sigma-1}$, there are $0<b_k\le B_k$ with $b_k(\sigma-1)^{-k^2}\le F(\sigma)\le B_k(\sigma-1)^{-k^2}$ for $1<\sigma\le2$.

\emph{Upper bound.} For $n\le x$, $\frac1n\le e\cdot n^{-1-1/\log x}$, so the sum is at most $eF(1+\frac1{\log x})\le eB_k(\log x)^{k^2}$.

\emph{Lower bound} (Rankin's trick backwards). Let $a\ge2$ and $\sigma_a=1+\frac a{\log x}$, assuming $\log x\ge a$. Then
\[
\sum_{n\le x}\frac{d_k(n)^2}n\ge\sum_{n\le x}\frac{d_k(n)^2}{n^{\sigma_a}}\ge F(\sigma_a)-\sum_{n>x}\frac{d_k(n)^2}{n^{\sigma_a}}\ge F(\sigma_a)-x^{-\frac{a-1}{\log x}}F\Bigl(1+\frac1{\log x}\Bigr),
\]
since $(n/x)^{(a-1)/\log x}\ge1$ for $n>x$. This is at least $(\log x)^{k^2}(b_ka^{-k^2}-e^{1-a}B_k)$, which is $\ge\frac12b_ka^{-k^2}(\log x)^{k^2}$ once $a$ is large enough in terms of $k$. For the finitely many smaller ranges of $x$, the sum is at least $1$.
\end{proof}

\begin{theorem}[Ramachandra; Heath-Brown; Rudnick and Soundararajan]\label{thm:momentlower}
For each integer $k\ge1$ there are $c_k>0$ and $T_k$ such that for $T\ge T_k$,
\[
\int_T^{2T}|\zeta(\tfrac12+it)|^{2k}\dd t\ge c_kT(\log T)^{k^2}.
\]
\end{theorem}
\begin{proof}
Let $x=T^{1/(8k)}$, $A(s)=\sum_{n\le x}n^{-s}$, and write $A(s)^j=\sum_nd_j(n;x)n^{-s}$, where $d_j(n;x)$ counts the $j$-tuples of integers $\le x$ with product $n$. So $0\le d_j(n;x)\le d_j(n)$, with equality for $n\le x$, and $\sum_nd_j(n;x)=\lfloor x\rfloor^j$. Put $\beta_m=d_{k-1}(m;x)$ ($\beta=\delta_1$ if $k=1$), $\gamma_n=d_k(n;x)$, and on the line $s=\frac12+it$,
\[
S_1=\int_T^{2T}\zeta(s)A(s)^{k-1}\overline{A(s)}^k\dd t,\qquad S_2=\int_T^{2T}|A(s)|^{2k}\dd t.
\]
By Hölder's inequality with exponents $2k$ and $\frac{2k}{2k-1}$,
\[
|S_1|\le\int|\zeta||A|^{2k-1}\le\bigl(\int|\zeta|^{2k}\bigr)^{1/(2k)}S_2^{1-1/(2k)},
\]
that is
\[
\int_T^{2T}|\zeta(\tfrac12+it)|^{2k}\dd t\ge\frac{|S_1|^{2k}}{S_2^{2k-1}}.
\]
It suffices to show $S_2\le C T(\log T)^{k^2}$ and $|S_1|\ge cT(\log T)^{k^2}$, with $C$, $c$ depending only on $k$.

\emph{The upper bound for $S_2$.} Write
\[
A(s)^k=\sum_{n\le x^k}\gamma_nn^{-1/2-it}.
\]
By Lemma~\ref{lem:mvt2}, with $\sum_n\gamma_nn^{1/2}\le x^{k/2}x^k$,
\[
S_2\le T\sum_{n\le x^k}\frac{d_k(n)^2}n+2x^{3k}\le C_kT(\log x^k)^{k^2}+2T^{3/8},
\]
by Lemma~\ref{lem:dk2}.

\emph{The lower bound for $S_1$.} By Volume~13, Lemma~1.5 (with its parameter equal to $5$ and $N=\lceil(2T+6)^{5/4}\rceil$), $\zeta(s)=\sum_{j\le N}j^{-s}+R(s)$ with $|R(s)|\le\frac{N^{1/2}}T+3N^{-1/2}\le CT^{-3/8}$ on $[T,2T]$. As $|A(s)|\le\sum_{n\le x}n^{-1/2}\le2x^{1/2}$, the part of $S_1$ coming from $R$ is at most $CT^{-3/8}\cdot T(2x^{1/2})^{2k-1}\le C'T^{3/4}$. The rest is
\[
\sum_{m,n}\beta_m\gamma_n\sum_{j\le N}\frac1{\sqrt{jmn}}\int_T^{2T}\Bigl(\frac n{jm}\Bigr)^{it}\dd t.
\]
The terms with $jm=n$ give $T\sum_{m\mid n}\beta_m\gamma_n\frac1n$ (here $j=\frac nm\le x^k\le N$). For fixed $m,n$, the terms with $jm\ne n$: if $jm\notin[\frac n2,2n]$, then $|\log\frac n{jm}|\ge\log2$ and they add up to at most $\frac2{\log2}\sum_{j\le N}j^{-1/2}\le6N^{1/2}$; if $jm\in[\frac n2,2n]$, then $|\log\frac n{jm}|\ge\frac{|n-jm|}{2n}\ge\frac1{2n}$, each term is at most $4n$ and there are at most $3n$ of them. So the off-diagonal part of $S_1$ is at most $\sum_{m,n}\beta_m\gamma_n(6N^{1/2}+12n^2)\le x^{2k-1}(6N^{1/2}+12x^{2k})\le CT^{7/8}$. Therefore
\[
S_1=T\sum_n\frac{\gamma_n}n\sum_{m\mid n}\beta_m+O(T^{7/8})\ge T\sum_{n\le x}\frac{d_k(n)^2}n+O(T^{7/8}),
\]
because all terms are nonnegative and, for $n\le x$, $\gamma_n=d_k(n)$ and $\sum_{m\mid n}\beta_m=\sum_{m\mid n}d_{k-1}(m)=d_k(n)$. By Lemma~\ref{lem:dk2}, $S_1\ge c_k(8k)^{-k^2}T(\log T)^{k^2}+O(T^{7/8})$.
\end{proof}

The argument is a cousin of the mollifier of Section~\ref{ch:selbergclt}. There $M\approx\zeta^{-1}$ made $\zeta M\approx1$; here $A^{k-1}\overline A^k$ is chosen so that $\zeta A^{k-1}$ imitates $\zeta^k$ and the correlation with $\overline A^k$ is as large as possible. Note also what the proof does \emph{not} give: the constant $c_k$ is much smaller than the conjectured one, because the length $x=T^{1/(8k)}$ of $A$ is short and Hölder's inequality is far from equality (Exercise~\ref{ex:k1}).

\section{The conjecture, and where its two factors come from}

\begin{proposition}\label{prop:eulermoment}
Let $(X_p)$ be independent and uniform on the unit circle, and $k\ge1$ an integer. Then
\[
\mathbb E\prod_{p\le y}\Bigl|1-\frac{X_p}{\sqrt p}\Bigr|^{-2k}=\prod_{p\le y}\sum_{r\ge0}\binom{k+r-1}r^2\frac1{p^r}\sim a_k\bigl(e^\gamma\log y\bigr)^{k^2}\qquad(y\to\infty),
\]
where
\[
a_k=\prod_p\Bigl(1-\frac1p\Bigr)^{k^2}\sum_{r\ge0}\binom{k+r-1}r^2\frac1{p^r}.
\]
\end{proposition}
\begin{proof}
By independence the expectation is a product. For $|z|<1$, $(1-z)^{-k}=\sum_r\binom{k+r-1}rz^r$, so $|1-X_p/\sqrt p|^{-2k}=|\sum_r\binom{k+r-1}rX_p^rp^{-r/2}|^2$, and $\mathbb EX_p^r\overline{X_p}^{r'}=\delta_{rr'}$ gives the factor $\sum_r\binom{k+r-1}r^2p^{-r}$. Write it as $(1-\frac1p)^{-k^2}h_0(\frac1p)$ with the function $h_0$ of the proof of Lemma~\ref{lem:dk2}; since $h_0(u)=1+O(u^2)$, $\prod_ph_0(\frac1p)$ converges to $a_k>0$, and $\prod_{p\le y}(1-\frac1p)^{-k^2}\sim(e^\gamma\log y)^{k^2}$ by Mertens.
\end{proof}

The random Euler product of Section~\ref{ch:bohrjessen} thus has moments of the conjectured shape, with the factor $a_k$ and the power $k^2$, but with $\log y$ in place of $\log T$, and with no $g_k$. A product over primes cannot see the zeros; random matrices see nothing else. Steve Gonek, Chris Hughes and Jon Keating found in 2007 how to combine the two. They proved that $\zeta(\frac12+it)$ factors, approximately, as an Euler product over the primes up to $y$ times a product over the zeros near $t$, and conjectured that the two factors are independent, the second behaving like $\det(I-U)$ for a random matrix of size $N=\frac{\log T}{e^\gamma\log y}$. With Proposition~\ref{prop:eulermoment} and Volume~11, Theorem~11.10, the moment of the product is then
\[
a_k(e^\gamma\log y)^{k^2}\cdot\frac{G(k+1)^2}{G(2k+1)}\Bigl(\frac{\log T}{e^\gamma\log y}\Bigr)^{k^2}=a_k\frac{G(k+1)^2}{G(2k+1)}(\log T)^{k^2},
\]
and $y$ disappears. This is the conjecture.

\medskip\noindent\textbf{The moment conjecture} (Keating and Snaith, 2000). For every integer $k\ge1$, and in fact every real $k>-\frac12$,
\[
\frac1T\int_0^T|\zeta(\tfrac12+it)|^{2k}\dd t\sim a_k\frac{G(k+1)^2}{G(2k+1)}(\log T)^{k^2}.
\]

\medskip
For $k=1$ this is Theorem~\ref{thm:ingham} ($a_1=1$, $\frac{G(2)^2}{G(3)}=1$), and for $k=2$ Ingham's theorem ($a_2=\frac6{\pi^2}$, $\frac{G(3)^2}{G(5)}=\frac1{12}$, product $\frac1{2\pi^2}$). Writing $\frac{G(k+1)^2}{G(2k+1)}=\frac{g_k}{(k^2)!}$ gives $g_3=42$ and $g_4=24024$ (Exercise~\ref{ex:gk}). Conrey, Farmer, Keating, Rubinstein and Snaith later conjectured the full asymptotic expansion, a polynomial of degree $k^2$ in $\log\frac t{2\pi}$, and it is that polynomial, not the leading term, that the numerics confirm.

\begin{figure}[ht]
\centering\fig{ch08-ratios}
\caption{The ratio of $\mathcal M_k(T)$ to the conjectured leading term $a_k\frac{g_k}{(k^2)!}\bigl(\log\frac t{2\pi}\bigr)^{k^2}$ (averaged over $[T,2T]$), estimated from $4\cdot10^5$ random points ($2\cdot10^5$ at $T=10^6$). For $k=1$ the ratio is $1.11$ at $T=10^6$, and Theorem~\ref{thm:ingham} explains the excess: the term $2\gamma-1$ of the mean square, against $\log\frac t{2\pi}\approx12.4$. For $k=2$ it is $2.6$, for $k=3$ about $47$: the lower-order terms of the polynomial dominate. The ratios do decrease, as slowly as powers of $\log T$ allow.}
\label{fig:momentratios}
\end{figure}

\begin{beauty}
The conjecture is a product of two worlds that do not know each other. The factor $a_k$ is pure arithmetic: it comes from independent random phases $p^{-it}$, one for each prime, as in every chapter of this book so far. The factor $\frac{G(k+1)^2}{G(2k+1)}$ is pure geometry of the unitary group, computed in Volume~11 by integrating over matrices, without a single prime in sight. That their product should be the moment of one function, $\zeta$, and that the same product should reproduce Hardy--Littlewood, Ingham and the number $42$ obtained by a quite different route, is one of the strongest pieces of evidence that the zeros of $\zeta$ behave like eigenvalues.
\end{beauty}

\begin{remark}
These are quoted for culture; see the Sources. (a) Upper bounds: on the Riemann hypothesis, Soundararajan proved $\mathcal M_k(T)\ll_{k,\varepsilon}(\log T)^{k^2+\varepsilon}$ for every $k>0$ (2009), and Harper removed the $\varepsilon$ (2013), with an argument close in spirit to the random multiplicative functions of Chapter~\ref{ch:randommult}. Unconditionally, the sharp upper bound $\mathcal M_k(T)\ll(\log T)^{k^2}$ is known for $0\le k\le2$ (Heath-Brown for $k=\frac1n$; Heap, Radziwi\l\l{} and Soundararajan in general, 2019). (b) Lower bounds: Theorem~\ref{thm:momentlower} extends to every real $k>0$ (Heath-Brown for rational $k$, 1981; Radziwi\l\l{} and Soundararajan for $k\ge1$; Heap and Soundararajan for all $k>0$, 2022). (c) The fourth moment is known with a power saving: Heath-Brown (1979) proved $\int_0^T|\zeta(\frac12+it)|^4\dd t=TP_4(\log T)+O(T^{7/8+\varepsilon})$ with a quartic polynomial $P_4$ of leading coefficient $\frac1{2\pi^2}$.
\end{remark}

\section{Exercises}

\begin{exercise}[$\star$]\label{ex:ak}
Show that $a_1=1$ and, using $\sum_{r\ge0}(r+1)^2u^r=\frac{1+u}{(1-u)^3}$, that $a_2=\prod_p(1-p^{-2})=\frac6{\pi^2}$.
\end{exercise}

\begin{exercise}[$\star$]\label{ex:gk}
Compute $\frac{G(k+1)^2}{G(2k+1)}(k^2)!$ for $k=1,2,3$, and check $g_3=42$. Show that $\frac{G(k+1)^2}{G(2k+1)}=\prod_{j=0}^{k-1}\frac{j!}{(j+k)!}$.
\end{exercise}

\begin{exercise}[$\star\star$]\label{ex:lognormal}
Let $Z\sim\mathcal N(0,v)$ with $v=\frac12\log\log T$. (a) Show that $\mathbb Ee^{2kZ}=(\log T)^{k^2}$. (b) Show that under the probability with density $e^{2kZ}/\mathbb Ee^{2kZ}$, $Z$ is $\mathcal N(2kv,v)=\mathcal N(k\log\log T,v)$. (c) Deduce that, if Selberg's law held exactly in the tail, the $2k$-th moment would be carried by the $t$ where $|\zeta(\frac12+it)|\approx(\log T)^k$, a set of probability $(\log T)^{-k^2+o(1)}$. Compare with Figure~\ref{fig:momentshare}.
\end{exercise}

\begin{exercise}[$\star\star$]\label{ex:diagonal4}
Assume Ramanujan's estimate $\sum_{n\le x}d(n)^2\sim\frac1{\pi^2}x\log^3x$. (a) Show that $\sum_{n\le x}\frac{d(n)^2}n\sim\frac1{4\pi^2}\log^4x$. (b) A heuristic writes $\zeta(\frac12+it)^2$ as $\sum_{n\le t/2\pi}d(n)n^{-1/2-it}$ plus a second sum of the same length, coming from the functional equation and nearly orthogonal to the first. Show that this heuristic predicts Ingham's constant $\frac1{2\pi^2}$, and explain the factor $2$.
\end{exercise}

\begin{exercise}[$\star$]\label{ex:k1}
Follow the proof of Theorem~\ref{thm:momentlower} for $k=1$ (so $S_1=\int\zeta\overline A$ and $S_2=\int|A|^2$) with $x=T^\theta$, $0<\theta<\frac13$. Show that it gives $\int_T^{2T}|\zeta(\frac12+it)|^2\dd t\ge(\theta+o(1))T\log T$, and compare with Theorem~\ref{thm:ingham}. Which inequality in the proof loses the factor $\theta$?
\end{exercise}

\begin{exercise}[$\star\star\star$]\label{ex:realk}
Let $k>0$ be real, and $\binom{k+r-1}r=\frac{k(k+1)\cdots(k+r-1)}{r!}$. Show that Proposition~\ref{prop:eulermoment} holds for every real $k>0$, and compute $a_{1/2}$ to three decimals with a computer. (The case $k=\frac12$ is the first moment of $|\zeta|$, for which no asymptotic formula is known.)
\end{exercise}

\section{Hints}

\begin{hint}{ex:ak}
For $k=1$ every factor is $(1-\frac1p)\sum_rp^{-r}=1$. For $k=2$ the factor is $(1-u)^4\frac{1+u}{(1-u)^3}$ with $u=\frac1p$.
\end{hint}
\begin{hint}{ex:gk}
$G(n+1)=G(n)\Gamma(n)=G(n)(n-1)!$.
\end{hint}
\begin{hint}{ex:lognormal}
Complete the square in $e^{2kz-z^2/(2v)}$. For (c), $P(Z\ge k\log\log T)=e^{-k^2(\log\log T)^2/(2v)(1+o(1))}$.
\end{hint}
\begin{hint}{ex:diagonal4}
Partial summation for (a). In (b), the second sum contributes the same diagonal as the first.
\end{hint}
\begin{hint}{ex:k1}
$S_1\approx T\sum_{n\le x}\frac1n$ and $S_2\approx T\sum_{n\le x}\frac1n$; the loss is in Cauchy--Schwarz.
\end{hint}
\begin{hint}{ex:realk}
$(1-z)^{-k}=\sum_r\binom{k+r-1}rz^r$ for $|z|<1$ and real $k$, and the coefficient of $u$ in $(1-u)^{k^2}\sum_r\binom{k+r-1}r^2u^r$ is still $k^2-k^2=0$.
\end{hint}

\section{Solutions}

\begin{solution}{ex:ak}
\idea The factors of $a_k$ are generating functions of squared binomial coefficients.
\writeup For $k=1$, $\binom r r=1$ and $(1-\frac1p)\sum_rp^{-r}=1$, so $a_1=1$. For $k=2$, $\binom{r+1}r=r+1$, and differentiating $\sum u^{r+1}=\frac u{1-u}$ twice in the form $\sum(r+1)u^r=\frac1{(1-u)^2}$, then $\sum(r+1)^2u^r=\frac{\dd}{\dd u}\frac u{(1-u)^2}=\frac{1+u}{(1-u)^3}$. So the factor is $(1-u)^4\frac{1+u}{(1-u)^3}=(1-u)(1+u)=1-u^2$ with $u=\frac1p$, and $a_2=\prod_p(1-p^{-2})=\frac1{\zeta(2)}=\frac6{\pi^2}$.
\end{solution}

\begin{solution}{ex:gk}
\idea Barnes' $G$ is a product of factorials.
\writeup $G(2)=G(3)=1$, $G(4)=2$, $G(5)=12$, $G(7)=0!\,1!\,2!\,3!\,4!\,5!=34560$. So $k=1$: $\frac{G(2)^2}{G(3)}\cdot1!=1$; $k=2$: $\frac1{12}\cdot4!=2$; $k=3$: $\frac4{34560}\cdot9!=\frac{362880}{8640}=42$. In general $G(k+1)=\prod_{j<k}j!$ and $G(2k+1)=\prod_{j<2k}j!=\prod_{j<k}j!\prod_{j<k}(j+k)!$, so $\frac{G(k+1)^2}{G(2k+1)}=\prod_{j<k}\frac{j!}{(j+k)!}$.
\end{solution}

\begin{solution}{ex:lognormal}
\idea Exponential tilting of a Gaussian shifts its mean.
\writeup (a) $\mathbb Ee^{\lambda Z}=e^{\lambda^2v/2}$ with $\lambda=2k$: $e^{2k^2v}=e^{k^2\log\log T}=(\log T)^{k^2}$. (b) $e^{2kz}e^{-z^2/(2v)}=e^{2k^2v}e^{-(z-2kv)^2/(2v)}$, so the tilted density is that of $\mathcal N(2kv,v)$, and $2kv=k\log\log T$. (c) Under the tilted law, $Z=k\log\log T+O_P(\sqrt{\log\log T})$, so the moment comes from $\log|\zeta|=k\log\log T(1+o(1))$, that is $|\zeta|=(\log T)^{k+o(1)}$. Under the original law this has probability $\exp(-\frac{(k\log\log T)^2}{2v}(1+o(1)))=(\log T)^{-k^2+o(1)}$. At $T=10^6$, $\log T\approx13.8$ and $(\log T)^{-k^2}$ is $0.07$, $4\cdot10^{-5}$ and $6\cdot10^{-11}$ for $k=1,2,3$: the last is far below what $2\cdot10^5$ samples can see, consistent with Figure~\ref{fig:momentshare}, where the share of the top $0.1\%$ is still growing with $k$.
\end{solution}

\begin{solution}{ex:diagonal4}
\idea The diagonal of a mean square counts squared coefficients with weight $\frac1n$.
\writeup (a) With $D(u)=\sum_{n\le u}d(n)^2\sim\frac1{\pi^2}u\log^3u$, $\sum_{n\le x}\frac{d(n)^2}n=\frac{D(x)}x+\int_1^x\frac{D(u)}{u^2}\dd u\sim\frac1{\pi^2}\int_1^x\frac{\log^3u}u\dd u=\frac1{4\pi^2}\log^4x$. (b) By Lemma~\ref{lem:mvt2}-type orthogonality (and ignoring that the length $t/2\pi$ exceeds the range where the lemma is efficient, which is why this is only a heuristic), each sum has mean square about $\sum_{n\le T/2\pi}\frac{d(n)^2}n\sim\frac1{4\pi^2}\log^4T$; if the two sums are nearly orthogonal, the total is twice that, $\frac1{2\pi^2}\log^4T$, Ingham's constant. The factor $2$ is the functional equation: for $\zeta^2$, a degree-$2$ object, both halves of the approximate functional equation have full length, and each carries half of the moment. For $|\zeta|^{2k}$ with $k\ge3$ the analogous diagonal heuristic gives the wrong constant, and it is the failure of such heuristics that made $g_3$ a mystery.
\end{solution}

\begin{solution}{ex:k1}
\idea Hölder's inequality is sharp only when $|\zeta|$ is proportional to $|A|$.
\writeup For $k=1$, $\beta=\delta_1$ and $\gamma_n=1$ for $n\le x$. The part of $S_1$ coming from $R$ is $O(T^{5/8}x^{1/2})$, and the off-diagonal part is $O(x(N^{1/2}+x^2))=O(T^{\theta+5/8}+T^{3\theta})$; both are $o(T)$ because $\theta<\frac13$. The diagonal is $T\sum_{n\le x}\frac1n$. So $S_1=T\sum_{n\le x}\frac1n+o(T)$, and Lemma~\ref{lem:mvt2} gives $S_2=T\sum_{n\le x}\frac1n+O(x^3)$, again with an error $o(T)$. (With sharper mean value theorems the same holds for every $\theta<\frac12$.) Both are $(\theta+o(1))T\log T$, and $\int|\zeta|^2\ge\frac{|S_1|^2}{S_2}=(\theta+o(1))T\log T$. Theorem~\ref{thm:ingham} gives $(1+o(1))T\log T$. The loss is in Cauchy--Schwarz, $|\int\zeta\overline A|^2\le\int|\zeta|^2\int|A|^2$: $\zeta$ is not proportional to $A$, it has a component orthogonal to all short Dirichlet polynomials, which carries the remaining $(1-\theta)T\log T$.
\end{solution}

\begin{solution}{ex:realk}
\idea Nothing in the proof used that $k$ is an integer.
\writeup For real $k>0$ and $|z|<1$, $(1-z)^{-k}=\exp(k\sum_jz^j/j)=\sum_r\binom{k+r-1}rz^r$ with the given coefficients, all positive. The computation $\mathbb E|1-X_p/\sqrt p|^{-2k}=\sum_r\binom{k+r-1}r^2p^{-r}$ is unchanged. With $h_0(u)=(1-u)^{k^2}\sum_r\binom{k+r-1}r^2u^r$, the coefficient of $u$ is $\binom k1^2-k^2=0$ and $h_0$ is analytic on $|u|<1$ with $h_0(0)=1$, so $h_0(u)=1+O_k(u^2)$ on $[0,\frac12]$ and $h_0>0$ there; the product $a_k=\prod_ph_0(\frac1p)$ converges to a positive number, and Mertens gives the asymptotic as before. Numerically (truncating at $p\le10^6$, the tail being $1+O(10^{-7})$), $a_{1/2}\approx0.988$. Every factor is close to $1$: expanding, $h_0(u)=1-\frac{k^2(k-1)^2}{4}u^2+O(u^3)$, and for $k=\frac12$ the coefficient is $-\frac1{64}$, so the product is roughly $\exp(-\frac1{64}\sum_pp^{-2})\approx e^{-0.007}$, the rest coming from the higher terms at $p=2$ and $p=3$.
\end{solution}

\section*{Sources and further reading}

\begin{itemize}
\item \emph{The first two moments.} G.~H.~Hardy and J.~E.~Littlewood, ``Contributions to the theory of the Riemann zeta-function and the theory of the distribution of primes'', \emph{Acta Math.}\ 41 (1918), 119--196. A.~E.~Ingham, ``Mean-value theorems in the theory of the Riemann zeta-function'', \emph{Proc.\ London Math.\ Soc.}\ (2) 27 (1926), 273--300. D.~R.~Heath-Brown, ``The fourth power moment of the Riemann zeta function'', \emph{Proc.\ London Math.\ Soc.}\ (3) 38 (1979), 385--422.
\item \emph{The numbers 42 and 24024.} J.~B.~Conrey and A.~Ghosh, ``A conjecture for the sixth power moment of the Riemann zeta-function'', \emph{Int.\ Math.\ Res.\ Not.}\ 1998, 775--780. J.~B.~Conrey and S.~M.~Gonek, ``High moments of the Riemann zeta-function'', \emph{Duke Math.\ J.}\ 107 (2001), 577--604. The number $42$ is in D.~Adams, \emph{The Hitchhiker's Guide to the Galaxy}, Pan Books, 1979.
\item \emph{Random matrices.} J.~P.~Keating and N.~C.~Snaith, ``Random matrix theory and $\zeta(1/2+it)$'', \emph{Comm.\ Math.\ Phys.}\ 214 (2000), 57--89. J.~B.~Conrey, D.~W.~Farmer, J.~P.~Keating, M.~O.~Rubinstein and N.~C.~Snaith, ``Integral moments of $L$-functions'', \emph{Proc.\ London Math.\ Soc.}\ (3) 91 (2005), 33--104. S.~M.~Gonek, C.~P.~Hughes and J.~P.~Keating, ``A hybrid Euler--Hadamard product for the Riemann zeta function'', \emph{Duke Math.\ J.}\ 136 (2007), 507--549.
\item \emph{Lower bounds.} K.~Ramachandra, ``Some remarks on the mean value of the Riemann zeta function and other Dirichlet series'', \emph{Hardy-Ramanujan J.}\ 1 (1978), 1--15. D.~R.~Heath-Brown, ``Fractional moments of the Riemann zeta function'', \emph{J.\ London Math.\ Soc.}\ (2) 24 (1981), 65--78. Z.~Rudnick and K.~Soundararajan, ``Lower bounds for moments of $L$-functions'', \emph{Proc.\ Natl.\ Acad.\ Sci.\ USA} 102 (2005), 6837--6838, whose argument we follow. W.~Heap and K.~Soundararajan, ``Lower bounds for moments of zeta and $L$-functions revisited'', \emph{Mathematika} 68 (2022), 1--14.
\item \emph{Upper bounds.} K.~Soundararajan, ``Moments of the Riemann zeta function'', \emph{Ann.\ of Math.}\ 170 (2009), 981--993. A.~J.~Harper, ``Sharp conditional bounds for moments of the Riemann zeta function'', arXiv:1305.4618 (2013). W.~Heap, M.~Radziwi\l\l{} and K.~Soundararajan, ``Sharp upper bounds for fractional moments of the Riemann zeta function'', \emph{Q.\ J.\ Math.}\ 70 (2019), 1387--1396.
\end{itemize}

